(1) 度规与 \(t,\phi\) 无关，故对有质量粒子
\[
\mathcal E=f\dot t,\qquad \ell=r^2\dot\phi,\qquad
-1=-f\dot t^2+f^{-1}\dot r^2+r^2\dot\phi^2
\]
（点表示对固有时 \(\tau\) 求导）。整理为
\[
\dot r^2+V_{\rm eff}(r)=\mathcal E^2,\qquad
V_{\rm eff}=f(r)\left(1+\frac{\ell^2}{r^2}\right).
\]
圆轨道同时满足 \(\dot r=0\)、\(V_{\rm eff}'=0\)。由于
\(V_{\rm eff}'=2M/r^2-2\ell^2/r^3+6M\ell^2/r^4\)，得到
\[
\boxed{\ell^2=\frac{Mr^2}{r-3M}},\qquad
\boxed{\mathcal E^2=\frac{(r-2M)^2}{r(r-3M)}}.
\]
必须 \(r>3M\) 才有有限的实数角动量；\(r=3M\) 不是有质量粒子的圆轨道。利用 \(\dot\phi/\dot t=f\ell/(r^2\mathcal E)\) 得
\(\boxed{\Omega=d\phi/dt=\sqrt{M/r^3}}\)，但它不意味着强场动力学等同牛顿力学。

(2) 保持 \(\ell\) 为原圆轨道值，令 \(r(\tau)=r_0+\rho(\tau)\)。由径向测地线方程
\(\ddot r=-V_{\rm eff}'/2\)，线性化为
\(\ddot\rho+\kappa_\tau^2\rho=0\)，其中
\[
\boxed{\kappa_\tau^2=\frac12 V_{\rm eff}''(r_0)
=\frac{M(r_0-6M)}{r_0^3(r_0-3M)}}.
\]
圆轨道 \(\dot t=\mathcal E/f=(1-3M/r_0)^{-1/2}\)，故坐标时本征频率为
\[
\boxed{\kappa_t^2=\frac{M(r_0-6M)}{r_0^4}
=\Omega^2\left(1-\frac{6M}{r_0}\right)}.
\]
当 \(3M<r_0<6M\) 时扰动指数增长；\(r_0>6M\) 时稳定，\(r_{\rm ISCO}=6M\) 是边界（恰在边界线性频率为零，须用高阶势判断）。近圆轨道每次径向周期的方位角为 \(2\pi\Omega/\kappa_t\)，相对 \(2\pi\) 的超额为
\[
\boxed{\Delta\phi=2\pi\left[\left(1-\frac{6M}{r_0}\right)^{-1/2}-1\right]
=\frac{6\pi M}{r_0}+O\!\left(\frac{M^2}{r_0^2}\right)}.
\]
该线性近心点公式仅适用于 \(r_0>6M\) 的小偏心率轨道，不能外推至 ISCO。

(3) 对零测地线取仿射参数 \(\lambda\)，有
\(E_\gamma=f\,dt/d\lambda\)、\(L_\gamma=r^2d\phi/d\lambda\)，于是
\[
\left(\frac{dr}{d\lambda}\right)^2+V_\gamma=E_\gamma^2,\qquad
V_\gamma=\frac{L_\gamma^2}{r^2}\left(1-\frac{2M}{r}\right).
\]
解 \(V_\gamma'=0\) 得 \(\boxed{r_{\rm ph}=3M}\)，且
\(V_\gamma''(3M)=-2L_\gamma^2/(81M^4)<0\)，故光子球不稳定。
临界冲量参数 \(b=L_\gamma/E_\gamma\) 满足 \(b_c^2=r_{\rm ph}^2/f(r_{\rm ph})\)，故
\[
\boxed{b_c=3\sqrt3\,M,\qquad \sigma_{\rm capture}=\pi b_c^2=27\pi M^2}.
\]
截面公式限于无穷远入射的几何光学零测地线；有质量轨道的有效势另含常数项 \(f\)，不可混用。

(4) 当地静止正交标架测得的轨道速度为
\[
\boxed{v=\frac{r\Omega}{\sqrt f}=\sqrt{\frac{M}{r-2M}}},\qquad
\gamma=(1-v^2)^{-1/2}=\sqrt{\frac{r-2M}{r-3M}},
\quad \frac{dt}{d\tau}=\frac{1}{\sqrt{1-3M/r}}.
\]
单位静质量的圆轨道束缚能为 \(1-\mathcal E\)，在 ISCO
\(\mathcal E_{\rm ISCO}=\sqrt{8/9}\)、\(\ell_{\rm ISCO}=2\sqrt3 M\)，
故理想薄盘常用的能量效率为 \(1-\sqrt{8/9}\)（不包含辐射俘获等修正）。
在发射者瞬时惯性系顺/逆切向发出的光子，经局域 Lorentz 变换到静止标架，能量各乘 \(\gamma(1\pm v)\)，再乘引力红移 \(\sqrt f\) 传至无穷远：
\[
\boxed{\frac{\nu_\infty^\pm}{\nu_0}
=\sqrt f\,\gamma(1\pm v)
=\frac{1-2M/r}{\sqrt{1-3M/r}}
\left(1\pm\sqrt{\frac{M}{r-2M}}\right)}.
\]
在 \(r=6M\) 时 \(v=1/2\)、\(\sqrt f\,\gamma=2\sqrt2/3\)，因此
\(\boxed{\nu_\infty^+/\nu_0=\sqrt2,\quad \nu_\infty^-/\nu_0=\sqrt2/3}\)。
两方向的切向发射光线在 \(r>3M\) 可向外逃逸；在 \(r\to\infty\) 时两频率比均趋于 \(1\)。当 \(r\downarrow3M\) 时圆轨道粒子的局域速度趋光速、所需能量与时间膨胀发散，而且轨道不稳定，不能把该极限理解为稳定的发射源。所有 \(M,r,b,\ell\) 在几何单位下均具有长度量纲，\(\Omega,\kappa_t\) 具有逆长度量纲，频率比与 \(\mathcal E\) 无量纲。
